\section{Chapter~\ref{sec:recognition}. Recognition}

The speed of recognition, the structure of the first stages, the linear readout of the response, and learning from the continuity of time are measured in humans and monkeys; the reduction of the whole chain to two operations and learning from video are tested on the book's models, and the explanations of illusions range from well founded to disputed.

{\footnotesize
\setlength{\tabcolsep}{3pt}\renewcommand{\arraystretch}{1.15}
\begin{longtable}{>{\raggedright}p{0.5cm}>{\raggedright}p{4.0cm}>{\raggedright}p{5.6cm}>{\raggedright}p{2.3cm}>{\raggedright\arraybackslash}p{1.7cm}}
\toprule
No. & Claim & Experiment and what was measured & Source & Status \\
\midrule\endfirsthead
\toprule
No. & Claim & Experiment and what was measured & Source & Status \\
\midrule\endhead
1 & Pixel-by-pixel template matching under shift does no better than a coin; a pair of stages~\eqref{eq:scstage} recognizes the shape at any position & \texttt{models/\allowbreak recognition\_models.py}, a ``retina'' of $24$ points, $4000$ trials: template, $0.49$, $0.51$, $0.52$ at flipped-point fractions $0$, $0.1$, $0.2$; the $S$ and $C$ pair, $1.00$, $0.86$, $0.70$ & derived in the chapter & model \\
2 & An animal in a picture is recognized in $\sim150$\,ms, a single pass through about ten stages of $10$--$20$\,ms each & Humans, ``animal or not'' task on photographs shown for $20$\,ms: evoked potentials for the two answers diverge after $\sim150$\,ms. Monkeys: the latency of the first response grows from stage to stage (V1 earliest, then V2 and V4). The number of stages and ``zero or one spike per stage'' are inferred from these numbers & \cite{Thorpe1996,Schmolesky1998} & measured \\
3 & Stage $S$ is AND over parts, stage $C$ is the maximum over positions (Proposition~\ref{prop:conveyor}) & Cat, primary visual cortex: simple cells respond to an edge of a given orientation at a given place, complex cells to the same orientation over a larger region. Intracellular recording from complex cells: for many cells the response to two bars is close to the larger of the two responses, on average closest to a max operation. The maximum via mutual inhibition is Proposition~\ref{prop:wta}; division by the total activity, a review & \cite{HubelWiesel1962,Lampl2004,Carandini2012} & measured \\
4 & Further up the chain the combinations are larger and more complex, and the response depends less and less on position & Monkeys, stimuli simplified down to the ``critical feature'': cells of V4 and posterior inferotemporal cortex respond to complex combinations of shape, color, and texture with small fields; in anterior inferotemporal cortex the fields are larger. Alternating $S$ and $C$ in a model reproduces this picture & \cite{KobatakeTanaka1994,Riesenhuber1999} & measured \\
5 & Convolutional networks repeat this alternation and predict the responses of the upper stages best; the object is read out linearly from the rates of a group of neurons & A network trained for recognition, with no fitting to the brain: the top layer predicts the responses of monkey inferotemporal neurons, the middle layers the responses of V4. A linear classifier on the activity of $\sim100$ randomly chosen cells, in windows from $12.5$\,ms, recognizes object and category across positions and sizes & \cite{Fukushima1980,Yamins2014,Hung2005} & measured \\
6 & The Hebb rule with a trace~\eqref{eq:traceinvariance} learns invariance from unlabeled video if the trace is at least $\sim20$\,ms & The idea of slow features and the trace rule is theory. \texttt{models/\allowbreak recognition\_models.py}, two $C$ neurons, $16$ inputs, $10$ runs, a $0.3$\,s pause between objects. At $\tau_{\text{tr}}=5$\,ms the match with the shape is $0.52$ on average, at $10$\,ms $0.56$; at $20$, $50$, and $200$\,ms, $1.00$ in all runs (at $30$\,ms one run in ten errs) & \cite{Foldiak1991,WiskottSejnowski2002} & model \\
7 & In the brain, invariance is learned from the continuity of time & Monkeys: the object was covertly swapped during an eye jump to a particular place in the field; the position invariance of inferotemporal neurons changed in line with the swap, significantly after just an hour. That this is the main rule of visual learning is a hypothesis & \cite{LiDiCarlo2008} & measured \\
8 & Motion: direction detectors, then the same AND/OR pair over large regions & Direction detectors in the rabbit retina; in monkey area MT all $110$ recorded neurons are direction selective, and the response to motion is stronger than in V1 & \cite{Barlow1965,Albright1984} & measured \\
9 & The upper stages send a prediction down, and the error goes up & The model explains extra-classical receptive field effects in primary visual cortex; individual observations agree (review), but it is not proven as the general design of the pipeline & \cite{RaoBallard1999,Keller2018} & hypothesis \\
10 & Hard pictures require feedback and several rounds & Monkeys: for ``hard'' pictures the decision in inferotemporal cortex forms $\sim30$\,ms later; these late responses are poorly predicted by a feedforward network and better by networks with feedback or by deeper ones & \cite{Kar2019} & measured \\
11 & An invisible pixel forgery fools a network; human vision is far more robust, but not immune & In networks: a small, specially chosen perturbation changes the answer; the ``panda $\to$ gibbon'' example is from the second paper. In humans, with brief presentation, forgeries that transfer well between networks shift the choice & \cite{Szegedy2014,Goodfellow2015,Elsayed2018} & measured \\
12 & Illusions of expectation: an unreliable input yields to expectation; faint things move ``more slowly,'' the dress is seen differently (Section~\ref{sec:illusions}) & Lower-contrast gratings appear to move more slowly. A survey of $1401$ people: the dress is seen in three stable variants, and a lighting cue switches the color; in $\sim13\,000$ observers the assumption about lighting decides. A Bayesian model with an expectation of slow motion explains a range of motion illusions & \cite{Thompson1982,LaferSousa2015,Wallisch2017,Weiss2002,Purves2011} & hypothesis \\
13 & Gain control: the waterfall, tilt, and contrast; the Hermann grid is not inhibition from neighbors & Direction detectors in the rabbit retina lower their rate under prolonged motion and fall below baseline after it stops. The tilt illusion is reproduced by a model with division by the activity of neighbors. Grid modifications that do not change the response of retinal output neurons abolish the spots, so the neighbor-inhibition explanation is rejected & \cite{BarlowHill1963,Schwartz2009,SchillerCarvey2005} & measured \\
14 & Delay compensation: a moving object appears ahead of a flash & The illusion is measured. Psychophysics contradicts both forward extrapolation of motion and a latency difference; a retrospective explanation has been proposed, based on events $\sim80$\,ms after the flash. The debate is not settled & \cite{Nijhawan1994,EaglemanSejnowski2000} & hypothesis \\
15 & Still ``snakes'' rotate: the latency difference~\eqref{eq:latency} is read by direction detectors & The illusion works without color too; pairs of neighboring elements produce it on their own; direction-selective neurons in monkey cortex give a directional response to these still pairs. In humans, rotation is triggered by microsaccades and blinks & \cite{Conway2005,OteroMillan2012} & measured \\
16 & Ambiguous pictures: mutual inhibition, fatigue, and noise; switches are caused mainly by noise & A model of competing neuron groups: switches in it are caused by \emph{noise}, and without noise there are none; it explains why the switching rate grows with stimulus strength. Fatigue plays a smaller role here & \cite{MorenoBote2007} & hypothesis \\
17 & Some illusions follow from the task the machine learns & A network trained to predict the next video frame ``sees'' rotation of the still rings and does not see it in control pictures; networks for image denoising and sharpening fall for color and brightness illusions, as humans do & \cite{Watanabe2018,GomezVilla2020} & measured \\
\bottomrule
\end{longtable}}
